% ============================================================================
% Chapter 7: Multi-Robot Swarm, Collaborative Mapping and Perception
% (Merged from former Ch.7 — Swarm/Communication and Ch.8 — Perception)
% ============================================================================

\chapter{Multi-Robot Swarm, Collaborative Mapping and Perception}
\label{chap:swarm}

% [TODO: short chapter lead-in: scaling from single-robot autonomy to a coordinated
%        swarm, and extending the stack with the perception pipeline for victim
%        detection in the SAR mission]

% ============================================================================
\section{Namespace-Based Scalability}
\label{sec:namespace_scalability}
% [TODO: ROS 2 namespacing of nodes/topics/TF; robot_name sets namespace AND
%        TF prefix; zero-code-change scaling — adding a third robot requires only
%        a new namespace value, no source changes]

\subsection{Firmware Namespace Support}
\label{subsec:firmware_namespace}
% [TODO: ESP32 receives namespace over serial before micro-ROS init;
%        micro_ros_agent namespaced via __ns; ensures joint_states and
%        joint_group_velocity_controller/commands are robot-specific]

\subsection{Namespaced Launch Architecture}
\label{subsec:namespaced_launch}
% [TODO: hardware_full_stack.launch.py staged bring-up; PushRosNamespace,
%        RewrittenYaml frame injection; per-robot TF frames (ausra_1/base_link,
%        ausra_2/base_link, etc.)]

% ============================================================================
\section{Replicating the System Across Robots}
\label{sec:replicating_robots}
% [TODO: copying the stack to robot 2 and robot 3; verification of correct
%        movement and mapping per robot; network configuration for multi-agent
%        ROS 2 discovery (ROS_DOMAIN_ID, FastDDS settings)]

% ============================================================================
\section{Inter-Robot Communication}
\label{sec:inter_robot_comms}
% [TODO: Wi-Fi-based communication; shared topics required for map merging
%        (/ausra_1/map, /ausra_2/map -> /map_merged); communication architecture;
%        bandwidth considerations for map topic relay]

% ============================================================================
\section{Collaborative Map Merging}
\label{sec:collaborative_map_merge}

\subsection{Map Expansion to a Fixed Global Canvas}
\label{subsec:map_expansion}
% [TODO: ausra_map_merge_HW; dynamic SLAM grid -> fixed 1000x1000 canvas;
%        heartbeat publisher keeps the merge node alive; segfault avoidance
%        strategy for grid-size mismatches]

\subsection{Known Initial Poses and Spawn Offsets}
\label{subsec:spawn_offsets}
% [TODO: tape-measured spawn offsets baked into canvas; known_init_poses=true;
%        merge into shared global map -> /map_merged; rationale for known-pose
%        strategy over feature-based merging in sparse disaster-area grids
%        (see also Section~\ref{subsec:map_merging} in the literature review)]

% ============================================================================
\section{Global Frontier-Based Exploration}
\label{sec:global_frontiers}
% [TODO: exploration driven by the shared /map_merged rather than individual
%        local maps; each robot's explore_lite instance subscribes to the merged
%        costmap; validated in simulation, ready for hardware transfer;
%        inter-robot frontier deconfliction strategy]

% ============================================================================
\section{Perception and Victim Detection}
\label{sec:perception}
\label{chap:perception}   % secondary alias so any future \ref{chap:perception} resolves

% [TODO: section lead-in — role of the OAK-D Lite and the YOLO-based detection
%        pipeline in the SAR mission; how the merged global map gives victim
%        detections a world-frame position; tight coupling with the swarm layer
%        motivates placing this here rather than in a separate chapter]

\subsection{Perception Requirements and Approach}
\label{sec:perception_requirements}
% [TODO: victim-detection objective; OAK-D Lite role (onboard MyriadX VPU);
%        on-robot inference vs. host inference trade-off; integration point
%        with hardware_full_stack.launch.py]

\subsection{Dataset and Model Training}
\label{sec:perception_dataset}
% [TODO: dataset collection and labeling methodology; victim classes defined;
%        training environment (Google Colab / local GPU); augmentation strategy]

\subsection{Model Iterations and Selection}
\label{sec:perception_models}
% [TODO: the multiple model versions/trials evaluated; key metrics (mAP, FPS
%        on MyriadX, false-positive rate); selection rationale for the final
%        deployed model]

\subsection{Simulation Perception Model}
\label{sec:perception_sim}
% [TODO: Gazebo victim model / actor; detection in the simulated environment;
%        sim-specific model weights; RViz visualization of detections]

\subsection{Real-World Perception Model}
\label{sec:perception_real}
% [TODO: model created for real-world victim detection; deployment on OAK-D Lite
%        via DepthAI SDK; latency and accuracy on physical hardware]

\subsection{Integration with the Autonomous Stack}
\label{sec:perception_integration}
% [TODO: ausra_yolo ROS 2 node; detection -> victim pose stamped in world frame
%        using depth + TF; reporting to base station; integration into
%        hardware_full_stack launch pipeline]

% ============================================================================
\section{Summary}
\label{sec:swarm_summary}
% [TODO: chapter summary covering both the multi-robot coordination layer
%        (namespacing, map merging, global frontier exploration) and the
%        perception pipeline (YOLO-based victim detection, OAK-D integration);
%        forward-pointer to Chapter~\ref{chap:results} for measured outcomes]
